\chapter{Free Molecules and Canonical Normal Forms}\label{ch:free}

A molecule is a composition of atoms, but a behavior admits many different
molecules, and what matters for engineering is structure: which atoms are
fused, in what order they run, how their wires are crossed. This chapter
separates the syntax of composition from its semantics. We construct
$F (\Sigma)$, the \emph{free symmetric monoidal category} on an atom signature
$\Sigma$: molecules built from the atoms of $\Sigma$, quotiented only by the
laws of a symmetric monoidal category. Up to coherence, a molecule is a
morphism of $F (\Sigma)$(a wiring diagram whose nodes are atoms) and its
meaning in Mol is the image of that diagram under the unique monoidal
evaluation.

Three facts about $F (\Sigma)$ do the work of the chapter. \emph{Existence and
the universal property:} every molecule over $\Sigma$ factors uniquely, up to
coherent isomorphism, through $F (\Sigma)$ (\autoref{thm:20}, \autoref{thm:21}). \emph{Finiteness:} for
a finite signature with bounded types and a bounded number of atoms, hom-sets
are finite (\autoref{thm:22}). \emph{Normal forms:} the runtime rewrite system of
\autoref{ch:rewrite} is terminating and confluent, so every molecule has a
unique normal form, and $M \simeq M'$ exactly when $\mathit{NF} (M) =
\mathit{NF} (M')$ (\autoref{thm:23}, \autoref{thm:24}). The pay-off is the \emph{precomputed lookup
table} of optimized molecules: one canonical, cost-optimal molecule per
behavior class, built once by exhaustive normalization (verified by the
\texttt{normal\_forms} tool, \autoref{app:a}) and consulted at runtime for
dense-integer dispatch; \autoref{sec:free-table} closes the chapter with it
and its role in the design problems of \autoref{ch:situations}.

\section{Types, atoms, and signatures}\label{sec:free-types}

\begin{definition}[Types and atom signatures]\label{def:types}
Fix a set $T$ of \emph{base types} (bytes, packets, contexts). A \emph{type}
is a finite sequence $A = (X_1,\dots,X_m)$ of base types, of \emph{length}
$|A| = m$; the empty sequence $I = ()$ is the unit type and the tensor is
concatenation, $A \otimes B = (X_1,\dots,X_m,Y_1,\dots,Y_n)$. A type universe
is \emph{bounded} when $T$ is finite and lengths are bounded by $d$; there
are then exactly $\sum_{m=0}^{d}|T|^m$ types. An \emph{atom signature}
$\Sigma$ is a set of \emph{atoms} typed $a \colon A \to B$ over the base
types. A \emph{molecule over} $\Sigma$ is an expression built from the atoms,
identities, and coherence isomorphisms by $\circ$ and $\otimes$.
\end{definition}

Types are exactly the objects of Mol, with tensor the tuple product and unit
the empty tuple (\autoref{ch:mol}, \autoref{thm:3}); Mol is a \emph{strict} symmetric
monoidal category, since tensor is literal concatenation. The behavioral
equivalence $\simeq$ (\autoref{thm:6}) is a congruence for $\circ$ and $\otimes$ and is
the relation for which the normal forms of this chapter are complete (\autoref{thm:24}).

\section{The free category on a signature}\label{sec:free-construction}

$F (\Sigma)$ is the quotient of the molecules over $\Sigma$ by the symmetric
monoidal axioms; we use the classical concrete model, morphisms as
\emph{wiring diagrams}, in which the axioms hold by construction.

\begin{definition}[The category $F (\Sigma)$]\label{def:fsigma}
$F (\Sigma)$ has objects the types over $\Sigma$ and morphisms $A \to B$ the
isomorphism classes of \emph{wiring diagrams} with boundary $(A,B)$: a finite
set of \emph{nodes} labeled by atoms $a \colon A_a \to B_a$ (input ports
labeled by the base types of $A_a$, output ports by those of $B_a$), the
$|A|$ boundary inputs and $|B|$ boundary outputs, and a \emph{wiring}: a
bijection between the source ports (boundary inputs and node outputs) and the
target ports (node inputs and boundary outputs) pairing only ports of the
same base type, such that the directed graph of nodes and wires is acyclic.
Composition $g \circ f$ plugs the boundary outputs of $f$ into the boundary
inputs of $g$; the identity wires each input to the corresponding output; the
tensor is disjoint union; the braiding $\sigma_{A,B} \colon A \otimes B \to
B \otimes A$ crosses the wires of $A$ and $B$.
\end{definition}

\begin{remark}
$F (\Sigma)$ is a small strict symmetric monoidal category: tensor is
concatenation, the unit is $I$, associator and unitors are identities, and
the braiding satisfies $\sigma_{B,A}\circ\sigma_{A,B} = \mathrm{id}$ and the
hexagon laws, so Mac Lane's coherence conditions hold \cite{maclane}.
Acyclicity is automatic for diagrams built from atoms (plugging never creates
a wire back to a lower layer) and is what separates the free category from
the traced structure of \autoref{ch:trace}.
\end{remark}

\section{Existence and the universal property}\label{sec:free-existence}

\begin{theorem}[Existence of $F (\Sigma)$]\label{thm:20}
 The free symmetric monoidal category $F (\Sigma)$ on an
atom signature $\Sigma$ exists: it is the initial algebra for the theory of
symmetric monoidal categories over $\Sigma$: the category of wiring
diagrams of \autoref{sec:free-construction}: and its existence also follows
abstractly from accessible and locally presentable categories \cite{adamek}.
\end{theorem}

\begin{proof}
\emph{Explicit construction.} The wiring diagrams form a symmetric monoidal
category: plugging is associative and unital; disjoint union is a bifunctor,
associative and unital with the empty diagram; the braiding is a natural
isomorphism satisfying the symmetry and hexagon laws. The category carries
the canonical $\Sigma$-structure $\rho_\Sigma$ (base type $X \mapsto (X)$;
atom $a \mapsto$ the one-node diagram) and is freely generated by it: every
morphism is a finite composite and tensor of atoms, identities, and braidings,
with no equations beyond the symmetric monoidal axioms, since the quotient by
them is part of the construction. Being initial among $\Sigma$-pointed
symmetric monoidal categories (\autoref{thm:21}), it is the free object.
\emph{Abstract existence.} Small symmetric monoidal categories are exactly
the models of a finitary essentially algebraic theory (category axioms,
tensor bifunctor, coherence isos and equations); such model categories are
locally presentable \cite{adamek}. The forgetful functor
$U \colon \mathbf{SMC} \to \mathbf{Graph}$ creates limits and preserves
filtered colimits (computed pointwise in $\mathbf{Graph}$), hence is
accessible, so by the adjoint functor theorem for locally presentable
categories: a limit-preserving accessible functor between such categories
has a left adjoint \cite{adamek}: $U$ has a left adjoint; applied to the
graph with vertices the base types and edges the atoms of $\Sigma$ it yields
$F (\Sigma)$. The two constructions coincide by the uniqueness in \autoref{thm:21}.
\end{proof}

\begin{theorem}[Universal property of $F (\Sigma)$]\label{thm:21}
 Let $\mathcal{C}$ be a symmetric monoidal category with a
$\Sigma$-structure $\rho$: an object $\rho (X)$ per base type $X$ and, per
atom $a \colon A \to B$, a morphism $\rho (a) \colon \rho (A) \to \rho (B)$
(with $\rho (A)$ the tensor of the images of the components of $A$). Then
there is a unique \emph{strict} monoidal functor $\Phi \colon F (\Sigma) \to
\mathcal{C}$ extending $\rho$; any strong monoidal functor extending $\rho$
is isomorphic to $\Phi$ by a unique monoidal natural isomorphism
\cite{maclane}. Consequently every molecule over $\Sigma$ factors uniquely,
up to coherent isomorphism, through $F (\Sigma)$.
\end{theorem}

\begin{proof}
\emph{Existence.} Set $\Phi (I) = I_{\mathcal{C}}$ and
$\Phi ((X_1,\dots,X_m)) = \rho (X_1) \otimes \cdots \otimes \rho (X_m)$; send
the one-node diagram of $a$ to $\rho (a)$, identities to identities, crossings
to the braiding of $\mathcal{C}$. Since a diagram is a finite arrangement of
atoms, identities, and braidings by $\circ$ and $\otimes$, and the evaluation
of a diagram in a symmetric monoidal category is well-defined by Mac Lane's
coherence theorem \cite{maclane}, $\Phi$ extends to all diagrams and is a
strict monoidal functor.
\emph{Uniqueness.} Any strict monoidal functor extending $\rho$ is forced on
objects (tensor of the $\rho (X)$), on one-node diagrams ($\rho (a)$), on
identities and crossings, and then on all diagrams by induction, as it
preserves $\circ$ and $\otimes$.
\emph{Up to coherent iso.} If $\Psi$ is strong monoidal, its constraints
$\varphi_{A,B} \colon \Psi (A) \otimes \Psi (B) \to \Psi (A \otimes B)$ define
$\alpha_{(X)} = \mathrm{id}$ and $\alpha_{A \otimes B} = \varphi_{A,B} \circ
(\alpha_A \otimes \alpha_B)$ by induction on the type length: the unique
monoidal natural isomorphism $\alpha \colon \Phi \Rightarrow \Psi$ (Mac
Lane's coherence theorem) \cite{maclane}.
\emph{Factorization.} Mol is strict symmetric monoidal (\autoref{thm:3}), so the atoms of
$\Sigma$ form a $\Sigma$-structure in Mol and $\mathrm{ev} \colon F (\Sigma)
\to \mathrm{Mol}$ is the unique strict monoidal functor extending it. A
molecule $M$ over $\Sigma$ is an expression, hence a unique morphism $[M]$ of
$F (\Sigma)$ up to coherence, with meaning $\mathrm{ev} ([M])$ in Mol; so
molecules factor through $F (\Sigma)$ as $\mathrm{ev}$, uniquely up to the
coherent isomorphism above.
\end{proof}

\section{Finiteness of hom-sets}\label{sec:free-finite}

\begin{theorem}[Finiteness for bounded signatures]\label{thm:22}
 Let $\Sigma$ be a finite atom signature with bounded
types: finitely many base types, types of length at most $d$, and every atom
with source and target of length at most $d$. Fix a bound $n$ on the number
of atoms in a molecule. Then for all types $A, B$ the set of morphisms of
$F (\Sigma)$ from $A$ to $B$ whose diagrams contain at most $n$ atoms is
finite. In particular, with a fixed atom bound only finitely many normal
forms exist (\autoref{thm:23}), so the lookup table of \autoref{thm:24} is finite and precomputable.
\end{theorem}

\begin{proof}
Only finitely many types occur, $\sum_{m=0}^{d}|T|^m$. A diagram $A \to B$
with $m \le n$ nodes has at most $s^m \le s^n$ labelings, $s$ being the
number of atoms. Each node has at most $d$ input and $d$ output ports,
so the diagram has at most $|A| + md \le d + nd$ source ports and at most
$|B| + md \le d + nd$ target ports, and a wiring is a label-respecting
bijection between them, of which there are at most $(d + nd)!$. Hence at most
$s^n (d + nd)!$ concrete diagrams, and no more isomorphism classes:
$F (\Sigma)(A,B)$ restricted to at most $n$ atoms is finite. No rewrite rule
introduces atoms (\autoref{thm:23}), so the normal forms of molecules with at most $n$
atoms are a subset of this finite set; by \autoref{thm:24} they are in bijection with the
behavior classes.
\end{proof}

\begin{remark}[The atom bound is essential]\label{rem:thm22-bound}
Bounded types alone do not make hom-sets finite: for a single base type $X$
and a single atom $f \colon X \to X$, the morphisms $\mathrm{id}, f,
f \circ f, \dots$ in $F (\Sigma)(X,X)$ are pairwise distinct, since the
universal property (\autoref{thm:21}) applied to the model in $\mathbf{Set}$ with
$X = \mathbb{N}$ and $f$ the successor function would otherwise force
$\mathrm{succ}^m = \mathrm{succ}^n$ for $m \neq n$. What a dataplane uses is
finiteness of \emph{bounded-size} molecules; a cost budget bounds the atom
count because every atom has positive cost (\autoref{thm:51}). This is the bounded,
enumerable class that the \texttt{normal\_forms} tool exhausts
(\autoref{app:a}).
\end{remark}

\section{Unique normal forms up to coherence}\label{sec:free-normal}

The free category fixes the \emph{structural} equations: molecules differing
only by coherence denote the same morphism of $F (\Sigma)$. The
\emph{behavioral} equations are handled by the runtime rewrite system of
\autoref{ch:rewrite}: the behavior-preserving rules of the runtime theories
ring cancellation (\autoref{thm:12}), parser left factoring (\autoref{thm:16}), batch flattening
(\autoref{thm:17}): read as local rewrites of a sub-diagram with fixed boundary, so
the system is a well-defined relation on the morphisms of $F (\Sigma)$.

\begin{theorem}[Unique normal forms]\label{thm:23}
 Let $R$ be the runtime rewrite system of
\autoref{ch:rewrite} on the molecules over a signature $\Sigma$. Then every
molecule $M$ over $\Sigma$ has a unique normal form $\mathit{NF} (M)$: the
unique $R$-irreducible morphism of $F (\Sigma)$ reachable from $M$, independent
of the order of rewriting. The normal form is unique up to coherence:
$\mathit{NF} (M)$ is a well-defined morphism of $F (\Sigma)$, and molecules
with the same wiring diagram have the same normal form. Moreover
$\mathit{NF} (M)$ contains at most as many atoms as $M$.
\end{theorem}

\begin{proof}
\emph{Termination.} No rule introduces atoms, and each rule strictly
decreases a well-founded measure: ring cancellation removes two nodes, batch
flattening merges nested batch nodes, left factoring lowers the measure of
\autoref{thm:16}; the measures of \autoref{thm:13}, \autoref{thm:16}, \autoref{thm:17} rule out infinite reductions.
\emph{Confluence.} The system is locally confluent: every critical pair joins
(\autoref{thm:14}, \autoref{thm:16}, \autoref{thm:17}), and redexes of different theories do not overlap, their
patterns being labeled by disjoint atoms. Termination and local confluence
imply confluence by Newman's lemma \cite{newman}.
\emph{Uniqueness.} By termination every reduction from $M$ reaches an
irreducible morphism; if $M \to^{*} u$ and $M \to^{*} v$ with $u, v$
irreducible, confluence gives a common reduct $w$ of $u$ and $v$, and
irreducibility forces $u = w = v$. Hence exactly one normal form
$\mathit{NF} (M)$ exists.
\emph{Up to coherence.} The rules replace a sub-diagram by another with the
same boundary, so they act on wiring diagrams: molecules with the same
diagram have the same normal form, a morphism of $F (\Sigma)$, and the atom
count does not increase.
\end{proof}

At the level of expressions the structural rules $(x \circ y) \circ z \to
x \circ (y \circ z)$, $\mathrm{id} \circ x \to x$, $x \circ \mathrm{id} \to
x$ compute canonical representatives: they are terminating (the sum over
composition nodes of the atoms in their left subtree strictly decreases) and
locally confluent (the critical pairs $(\mathrm{id} \circ x) \circ z$,
$(x \circ y) \circ \mathrm{id}$, $(x \circ \mathrm{id}) \circ z$,
$((x \circ y) \circ z) \circ w$ join in at most three steps), hence confluent
\cite{newman}, with normal forms the right-nested chains without identities;
they preserve the wiring diagram, so they may be interleaved with the
behavioral rules without changing $\mathit{NF} (M)$.

\begin{theorem}[Completeness and the lookup table]\label{thm:24}
 Let $R$ be the runtime rewrite system of
\autoref{ch:rewrite} on the molecules over a signature $\Sigma$: terminating
and confluent (\autoref{thm:13}--\autoref{thm:17}) and complete for behavioral equivalence, in that
the congruence generated by the rules of $R$ is exactly $\simeq$ on the
molecules over $\Sigma$ (a complete system in the sense of Knuth--Bendix
\cite{knuthbendix}; established for the runtime theories by \autoref{thm:15}, \autoref{thm:18},
\autoref{thm:19}). Then for all molecules $M, M'$ over $\Sigma$:
\[
M \simeq M' \quad\text{iff}\quad \mathit{NF} (M) = \mathit{NF} (M'),
\]
and $\mathit{NF} (M) \simeq M$ for every $M$. If $\Sigma$ is finite with
bounded types and bounded atom count, the normal forms form a finite,
precomputable lookup table of optimized molecules, one entry per behavior
class.
\end{theorem}

\begin{proof}
\emph{Soundness.} Each rule preserves behavior and is cost-nonincreasing
(\autoref{thm:12}, \autoref{thm:16}--\autoref{thm:18}), so $M \to^{*} \mathit{NF} (M)$ implies $M \simeq
\mathit{NF} (M)$ by induction on the reduction.
\emph{Completeness $(\Rightarrow)$.} By hypothesis $M \simeq M'$ means that
$M$ and $M'$ are connected by $R$-steps, and this relation is joinability:
joinability is reflexive, symmetric, transitive, and preserved by single
steps on either side (if $u \downarrow v$ and $u' \to u$ then $u' \to^{*}
w$; if $v \to v'$, confluence applied to $v$ yields a common reduct of $w$
and $v'$). Hence $M$ and $M'$ have a common reduct, which by termination and
confluence reduces to $\mathit{NF} (M) = \mathit{NF} (M')$.
\emph{Completeness $(\Leftarrow)$.} If $\mathit{NF} (M) = \mathit{NF} (M')$,
then $M \simeq \mathit{NF} (M) \simeq \mathit{NF} (M') \simeq M'$ by soundness.
\emph{Lookup table.} By \autoref{thm:22} the normal forms of bounded molecules are finite
in number, and the two completeness directions make $M \mapsto \mathit{NF} (M)$
a bijection between behavior classes and normal forms; hence the table is
one cost-optimal molecule per class (\autoref{thm:26}, \autoref{thm:29}), indexed by the type pair
and a dense integer: is finite and precomputed by exhaustive normalization.
The \texttt{normal\_forms} tool (\autoref{app:a}) performs exactly this
enumeration on the signature $f \colon A \to B$, $g \colon B \to C$,
$h \colon C \to D$: all $3{,}929$ well-typed terms with at most six atom
occurrences reduce to the $10$ distinct normal forms, with equality of normal
forms agreeing with equality of the realized functions on all
$7{,}716{,}556$ pairs: PASS.
\end{proof}

\section{The precomputed lookup table}\label{sec:free-table}

A dataplane cannot optimize at runtime; it dispatches. The lookup table is the
theory of this chapter compiled into a finite artifact: for the runtime
signature (finite, bounded types, bounded hot-path size) it assigns to each
type pair $A \to B$ its finitely many behavior classes (\autoref{thm:22}), each
represented by its normal form (\autoref{thm:23}), which is exactly its behavior (\autoref{thm:24});
with the costs of \autoref{ch:enrich} the stored representative is
cost-optimal within its class (\autoref{thm:26}, \autoref{thm:29}). Using the table is a lookup
followed by dense-integer dispatch (policy [MOL] of the companion standard):
constant-time and allocation-free.

For the signature $f \colon A \to B$, $g \colon B \to C$, $h \colon C \to D$
the table is the following; the third column counts the well-typed terms of
depth at most six that normalize to each entry (\texttt{normal\_forms},
\autoref{app:a}).

\begin{table}[htbp]
\centering
\begin{tabular}{@{}llr@{}}
\toprule
\textbf{Normal form $\mathit{NF} (M)$} & \textbf{Type} & \textbf{Terms in class} \\
\midrule
$\mathrm{id}_A$         & $A \to A$ & 65   \\
$f$                     & $A \to B$ & 351  \\
$g \circ f$             & $A \to C$ & 807  \\
$h \circ g \circ f$     & $A \to D$ & 1002 \\
$\mathrm{id}_B$         & $B \to B$ & 65   \\
$g$                     & $B \to C$ & 351  \\
$h \circ g$             & $B \to D$ & 807  \\
$\mathrm{id}_C$         & $C \to C$ & 65   \\
$h$                     & $C \to D$ & 351  \\
$\mathrm{id}_D$         & $D \to D$ & 65   \\
\midrule
Total                   &           & 3929 \\
\bottomrule
\end{tabular}
\caption{The lookup table for the signature $\{f, g, h\}$: the ten behavior
classes (normal forms) and, for each, the number of enumerated well-typed
terms of depth at most six that normalize to it.}
\label{tab:lookup}
\end{table}

Each row is the unique normal form of its behavior class (\autoref{thm:23}, \autoref{thm:24}), so the
table is consistent by construction; syntactically distinct terms in a class
one entry per class: $h \circ (g \circ f)$ and $(h \circ g) \circ f$, or $f$ and
$f \circ \mathrm{id}_A$: share one entry: the fusion promised by cut
elimination (\autoref{ch:linear}, \autoref{thm:49}) is, for bounded signatures, a table
lookup. The same enumeration resolves the design decision problems of
\autoref{ch:situations}: a situation $(A \to B, \Phi, \bar{c})$ seeks a
molecule realizing class $\Phi$ within cost $\bar{c}$, and the table restricts
the search to one candidate per class, computed once (\autoref{thm:51}).

\begin{remark}[Summary]
$F (\Sigma)$ is the syntax of molecules: existence and the universal property
(\autoref{thm:20}, \autoref{thm:21}) make it the canonical interface between expressions and Mol;
finiteness (\autoref{thm:22}) makes optimization exhaustive; unique normal forms (\autoref{thm:23})
and completeness (\autoref{thm:24}) make it unambiguous, so every behavior class has
exactly one canonical, cost-optimal representative, precomputable before
runtime and dispatched in constant time.
\end{remark}
